% Options for packages loaded elsewhere
% Options for packages loaded elsewhere
\PassOptionsToPackage{unicode}{hyperref}
\PassOptionsToPackage{hyphens}{url}
\PassOptionsToPackage{dvipsnames,svgnames,x11names}{xcolor}
%
\documentclass[
  18,
]{article}
\usepackage{xcolor}
\usepackage[top=10mm,left=15mm,right=15mm,bottom=5mm,heightrounded]{geometry}
\usepackage{amsmath,amssymb}
\setcounter{secnumdepth}{-\maxdimen} % remove section numbering
\usepackage{iftex}
\ifPDFTeX
  \usepackage[T1]{fontenc}
  \usepackage[utf8]{inputenc}
  \usepackage{textcomp} % provide euro and other symbols
\else % if luatex or xetex
  \usepackage{unicode-math} % this also loads fontspec
  \defaultfontfeatures{Scale=MatchLowercase}
  \defaultfontfeatures[\rmfamily]{Ligatures=TeX,Scale=1}
\fi
\usepackage[]{libertine}
\ifPDFTeX\else
  % xetex/luatex font selection
\fi
% Use upquote if available, for straight quotes in verbatim environments
\IfFileExists{upquote.sty}{\usepackage{upquote}}{}
\IfFileExists{microtype.sty}{% use microtype if available
  \usepackage[]{microtype}
  \UseMicrotypeSet[protrusion]{basicmath} % disable protrusion for tt fonts
}{}
\makeatletter
\@ifundefined{KOMAClassName}{% if non-KOMA class
  \IfFileExists{parskip.sty}{%
    \usepackage{parskip}
  }{% else
    \setlength{\parindent}{0pt}
    \setlength{\parskip}{6pt plus 2pt minus 1pt}}
}{% if KOMA class
  \KOMAoptions{parskip=half}}
\makeatother
% Make \paragraph and \subparagraph free-standing
\makeatletter
\ifx\paragraph\undefined\else
  \let\oldparagraph\paragraph
  \renewcommand{\paragraph}{
    \@ifstar
      \xxxParagraphStar
      \xxxParagraphNoStar
  }
  \newcommand{\xxxParagraphStar}[1]{\oldparagraph*{#1}\mbox{}}
  \newcommand{\xxxParagraphNoStar}[1]{\oldparagraph{#1}\mbox{}}
\fi
\ifx\subparagraph\undefined\else
  \let\oldsubparagraph\subparagraph
  \renewcommand{\subparagraph}{
    \@ifstar
      \xxxSubParagraphStar
      \xxxSubParagraphNoStar
  }
  \newcommand{\xxxSubParagraphStar}[1]{\oldsubparagraph*{#1}\mbox{}}
  \newcommand{\xxxSubParagraphNoStar}[1]{\oldsubparagraph{#1}\mbox{}}
\fi
\makeatother


\usepackage{longtable,booktabs,array}
\usepackage{calc} % for calculating minipage widths
% Correct order of tables after \paragraph or \subparagraph
\usepackage{etoolbox}
\makeatletter
\patchcmd\longtable{\par}{\if@noskipsec\mbox{}\fi\par}{}{}
\makeatother
% Allow footnotes in longtable head/foot
\IfFileExists{footnotehyper.sty}{\usepackage{footnotehyper}}{\usepackage{footnote}}
\makesavenoteenv{longtable}
\usepackage{graphicx}
\makeatletter
\newsavebox\pandoc@box
\newcommand*\pandocbounded[1]{% scales image to fit in text height/width
  \sbox\pandoc@box{#1}%
  \Gscale@div\@tempa{\textheight}{\dimexpr\ht\pandoc@box+\dp\pandoc@box\relax}%
  \Gscale@div\@tempb{\linewidth}{\wd\pandoc@box}%
  \ifdim\@tempb\p@<\@tempa\p@\let\@tempa\@tempb\fi% select the smaller of both
  \ifdim\@tempa\p@<\p@\scalebox{\@tempa}{\usebox\pandoc@box}%
  \else\usebox{\pandoc@box}%
  \fi%
}
% Set default figure placement to htbp
\def\fps@figure{htbp}
\makeatother





\setlength{\emergencystretch}{3em} % prevent overfull lines

\providecommand{\tightlist}{%
  \setlength{\itemsep}{0pt}\setlength{\parskip}{0pt}}



 


\usepackage{multicol}
\usepackage{lipsum}
\usepackage{hyperref}
\usepackage{libertine}
\usepackage{fancyhdr}
\usepackage{lastpage}
\setlength{\multicolsep}{2pt plus 1.0pt minus 0.75pt}
\setlength{\columnsep}{3em}
\usepackage{enumitem}
\setlist{nolistsep}
\setlength{\columnsep}{-2.8cm}
\fancyhf{}
\renewcommand{\headrulewidth}{0pt}
\fancyfoot[C]{\thepage/\pageref{LastPage}}
\AtBeginDocument{\pagestyle{fancy}}
\makeatletter
\@ifpackageloaded{caption}{}{\usepackage{caption}}
\AtBeginDocument{%
\ifdefined\contentsname
  \renewcommand*\contentsname{Table of contents}
\else
  \newcommand\contentsname{Table of contents}
\fi
\ifdefined\listfigurename
  \renewcommand*\listfigurename{List of Figures}
\else
  \newcommand\listfigurename{List of Figures}
\fi
\ifdefined\listtablename
  \renewcommand*\listtablename{List of Tables}
\else
  \newcommand\listtablename{List of Tables}
\fi
\ifdefined\figurename
  \renewcommand*\figurename{Figure}
\else
  \newcommand\figurename{Figure}
\fi
\ifdefined\tablename
  \renewcommand*\tablename{Table}
\else
  \newcommand\tablename{Table}
\fi
}
\@ifpackageloaded{float}{}{\usepackage{float}}
\floatstyle{ruled}
\@ifundefined{c@chapter}{\newfloat{codelisting}{h}{lop}}{\newfloat{codelisting}{h}{lop}[chapter]}
\floatname{codelisting}{Listing}
\newcommand*\listoflistings{\listof{codelisting}{List of Listings}}
\makeatother
\makeatletter
\makeatother
\makeatletter
\@ifpackageloaded{caption}{}{\usepackage{caption}}
\@ifpackageloaded{subcaption}{}{\usepackage{subcaption}}
\makeatother
\usepackage{bookmark}
\IfFileExists{xurl.sty}{\usepackage{xurl}}{} % add URL line breaks if available
\urlstyle{same}
\hypersetup{
  pdfauthor={Dr.~Simone Romiti},
  colorlinks=true,
  linkcolor={blue},
  filecolor={Maroon},
  citecolor={Blue},
  urlcolor={Blue},
  pdfcreator={LaTeX via pandoc}}


\author{}
\date{}
\begin{document}


\begin{huge}\begin{center}{\bf SIMONE ROMITI}\end{center}\end{huge}

\begin{center}
simone.romiti.1994@gmail.com |
\href{https://simone-romiti.github.io/}{Website} | 
\href{https://github.com/simone-romiti}{GitHub} | 
\href{https://www.linkedin.com/in/simone-romiti/}{LinkedIn} |
\href{https://orcid.org/0000-0002-6509-447X}{ORCID}
\end{center}
\vspace{15pt}

\hrule

\begin{large}



Theoretical physicist specialized in lattice gauge theories and lattice QCD predictions, with 8+ years of experience in high-performance computing, large-scale numerical simulations, and machine learning. 
Proven track record of designing algorithms that delivered order-of-magnitude performance gains, building open-source scientific software used in peer-reviewed research, and leading collaborative projects across international research institutions.




\end{large}

\vspace{9pt}

\begin{large}{\bf TECHNICAL SKILLS}
  \vspace{9pt}
  \hrule
  \vspace{9pt}
\end{large}
\vspace{-0.15cm}

\begin{large}

\textbf{Programming Languages} - C, C++, Python, Bash, R

\vspace{-0.15cm}
\textbf{HPC \& Systems} - CUDA, MPI, OpenMP, SLURM, EasyBuild, GNU/Linux

\vspace{-0.15cm}
\textbf{ML \& Data} - PINNs, VAEs, Diffusion Models, Bayesian Statistics, Markov Chain Monte Carlo, Nested Sampling

\vspace{-0.15cm}
\textbf{Libraries \& Frameworks} - PyTorch, NumPy, SciPy, Pandas, Matplotlib, Plotly, Streamlit, SymPy, Jupyter

\vspace{-0.15cm}
\textbf{DevOps \& Tools} - Git, GitHub Actions, CI/CD, Docker, \LaTeX, Quarto, Markdown

\vspace{-0.15cm}
\textbf{Spoken Languages} - Italian (native), English (fluent), German (basic)

\end{large}

\vspace{9pt}

\begin{large}{\bf WORK EXPERIENCE}
  \vspace{9pt}
  \hrule
  \vspace{9pt}

  \begin{multicols}{2}
    \begin{flushleft}{\bf \href{https://www.unibe.ch/index_eng.html}{University of Bern}}\end{flushleft}
    \begin{flushright}Apr 2024 -- Present\end{flushright}
  \end{multicols}
  \vspace{-0.15cm}
  \begin{multicols}{2}
    \begin{flushleft}\textit{Postdoctoral Research Scientist}\end{flushleft}
    \begin{flushright}Bern, Switzerland\end{flushright}
  \end{multicols}
\end{large}

\vspace{-0.09cm}

\begin{large}
  \begin{itemize}
    \item \textbf{Reduced memory complexity from exponential to polynomial} by designing a novel algorithm using Physics-Informed Neural Networks (PINNs) in PyTorch
    \item \textbf{Achieved sub-permille numerical precision} as lead developer for a high-accuracy scientific computation pipeline
    \item \textbf{Improved algorithm scaling from $O(N^6)$ to $O(N \log N)$} through a novel mathematical reformulation using FFTs and convolutions
    \item \textbf{Built and maintained open-source simulation libraries} (Python/C++) adopted across multiple research groups and directly enabling peer-reviewed publications
    \item \textbf{Mentored 3 PhD students} on software engineering practices, numerical methods, and project delivery
  \end{itemize}
\end{large}

\vspace{9pt}

\begin{large}
  \begin{multicols}{2}
    \begin{flushleft}{\bf \href{https://www.uni-bonn.de/en/home?set_language=en}{University of Bonn}}\end{flushleft}
    \begin{flushright}Nov 2021 -- Mar 2024\end{flushright}
  \end{multicols}
  \vspace{-0.15cm}
  \begin{multicols}{2}
    \begin{flushleft}\textit{Postdoctoral Research Scientist}\end{flushleft}
    \begin{flushright}Bonn, Germany\end{flushright}
  \end{multicols}
\end{large}

\vspace{-0.09cm}

\begin{large}
  \begin{itemize}
    \item \textbf{Optimized GPU kernels (CUDA)}, achieving a $\sim$1.5$\times$ throughput improvement via auto-tuning of Multigrid solver parameters on HPC clusters
    \item \textbf{Designed and ran large-scale distributed simulations} (MPI/OpenMP) on international computing centers
    \item \textbf{Developed a novel numerical method} achieving machine-precision exactness for a class of operator algebra constraints
    \item \textbf{Integrated Monte Carlo simulations with quantum computing workflows}, delivering end-to-end prototype pipelines bridging classical HPC and quantum backends
    \item \textbf{Mentored 3 graduate students} and delivered tutorial sessions for undergraduate computing courses
  \end{itemize}
\end{large}

\vspace{9pt}

\begin{large}{\bf SELECTED PUBLICATIONS}
  \vspace{9pt}
  \hrule
  \vspace{9pt}

   \begin{itemize}
     \item \href{https://inspirehep.net/literature/3075565}{SU(N) lattice gauge theories with Physics-Informed Neural Networks}
     \item \href{https://inspirehep.net/literature/2925594}{The anomalous magnetic moment of the muon in the Standard Model: an update}
     \item \href{https://inspirehep.net/literature/2847988}{Strange and charm quark contributions to the muon anomalous magnetic moment in lattice QCD with twisted-mass fermions}
     \item \href{https://inspirehep.net/literature/2781250}{Towards determining the (2+1)-dimensional Quantum Electrodynamics running coupling with Monte Carlo and quantum computing methods}
     \item \href{https://inspirehep.net/literature/2724218}{Digitizing lattice gauge theories in the magnetic basis: reducing the breaking of the fundamental commutation relations}
   \end{itemize}

\end{large}

\vspace{9pt}

\begin{large}{\bf EDUCATION}
  \vspace{9pt}
  \hrule
  \vspace{9pt}

  \begin{multicols}{2}
    \begin{flushleft}{\href{https://www.uniroma3.it/}{\textit{Roma Tre University}}}\end{flushleft}
    \begin{flushright}Rome, Italy\end{flushright}
  \end{multicols}
  \vspace{-0.15cm}
  \begin{multicols}{2}
    \begin{flushleft}
    \textbf{PhD in Physics} (Dissertation April 2022) 
    \end{flushleft}
    \begin{flushright}Nov. 2018 -- Oct. 2021\end{flushright}
  \end{multicols}
  \vspace{-0.2cm}
  \begin{itemize}
    \item Thesis: \textit{Non-perturbative aspects of gauge theories on the lattice: algorithms, simulations and quantum computing}
    \item Focus: Monte Carlo methods, HPC simulations, Isospin Breaking effects in lattice QCD
    \item Ranked \textbf{1st} in national competitive admission exam
  \end{itemize}

  \begin{multicols}{2}
    \begin{flushleft}{\href{https://www.uniroma3.it/}{\textit{Roma Tre University}}}\end{flushleft}
    \begin{flushright}Rome, Italy\end{flushright}
  \end{multicols}
  \vspace{-0.15cm}
  \begin{multicols}{2}
    \begin{flushleft}
    \textbf{M.S.\ in Physics}
    \end{flushleft}
    \begin{flushright}Oct. 2016 -- Oct. 2018\end{flushright}
  \end{multicols}
  \vspace{-0.2cm}
  \begin{itemize}
    \item GPA: 29.85/30, graduated with highest honours
  \end{itemize}

  \begin{multicols}{2}
    \begin{flushleft}{\href{https://www.uniroma3.it/}{\textit{Roma Tre University}}}\end{flushleft}
    \begin{flushright}Rome, Italy\end{flushright}
  \end{multicols}
  \vspace{-0.15cm}
  \begin{multicols}{2}
    \begin{flushleft}
    \textbf{B.S.\ in Physics}
    \end{flushleft}
    \begin{flushright}Oct. 2013 -- Jul. 2016\end{flushright}
  \end{multicols}
  \vspace{-0.2cm}
  \begin{itemize}
    \item GPA: 28.84/30, graduated with highest honours | Merit Scholarship recipient
  \end{itemize}

\end{large}

\vspace{9pt}

\begin{large}
{\bf LEADERSHIP \& ACHIEVEMENTS}
  \vspace{9pt}
  \hrule
  \vspace{9pt}

\setlength{\columnsep}{-4.5cm}

  \begin{multicols}{2}
    \begin{flushleft}    
      \textbf{Principal Investigator} -- 0.24M GPU node-hours compute grant
    \end{flushleft}
    \begin{flushright}\href{https://www.cscs.ch/computers/alps}{CSCS ALPS} | 2025\end{flushright}
  \end{multicols}

  \vspace{-0.15cm}

  \begin{multicols}{2}
    \begin{flushleft}    
      \textbf{Workshop Organizer} -- HPC \& Quantum Computing workshop
    \end{flushleft}
    \begin{flushright}\href{https://indico.ectstar.eu/event/242/}{ECT*} | 2025\end{flushright}
  \end{multicols}

  \vspace{-0.15cm}

  \begin{multicols}{2}
    \begin{flushleft}    
      \textbf{Invited Speaker} -- Scientific seminars at \href{https://indico.cern.ch/category/376/}{CERN} and \href{https://indico.ectstar.eu/event/229/contributions/5308/}{ECT*}
    \end{flushleft}
    \begin{flushright}2025 -- 2026\end{flushright}
  \end{multicols}

  \vspace{-0.15cm}  

  \begin{multicols}{2}
    \begin{flushleft}    
      \textbf{Visiting Research Grant} -- \href{https://www.philnat.unibe.ch/e17097/e95924/e1163513/e1163533/OpenRound2026_01_2_eng.pdf}{Open round 2026/1 fund}
    \end{flushleft}
    \begin{flushright}\href{https://www.tuwien.at/}{TU Wien} | 2026\end{flushright}
  \end{multicols}


\end{large}




\end{document}
